% preface_section1_v2.tex — Section 1 of the Preface (v2)
% Macros assumed from main.tex preamble
\providecommand{\cA}{\mathcal{A}}
\providecommand{\barB}{\overline{B}}
\providecommand{\gAmod}{\mathfrak{g}_{\cA}^{\mathrm{mod}}}
\providecommand{\MC}{\text{MC}}

\section*{1.\quad The bar construction on algebraic curves}
% ====================================================================

\medskip

\noindent
A chiral algebra on a smooth algebraic curve~$X$ encodes
operator product singularities: the product of two sections
$a(z)\,b(w)$ diverges as $z\to w$, and the divergence carries
algebraic data.  The problem is to extract that data systematically
at all orders, on all curves, at all genera.  The solution is
classical: the bar construction, applied not to an associative
algebra on a point but to a chiral algebra on a curve.  The integral
kernel is the logarithmic $1$-form $d\log(z_1-z_2)$; the
compactification is Fulton--MacPherson; the coalgebra is cofree
on desuspended generators; the differential squares to zero if and
only if the algebra satisfies the Borcherds identity.

% ====================================================================
\subsection*{The propagator}
% ====================================================================

\noindent
On the configuration space $C_2(X) = X \times X \setminus \Delta$
of a smooth algebraic curve~$X$ over~$\mathbf{C}$, define
\begin{equation}\label{eq:pf1-eta}
\eta_{12} \;=\; d\log(z_1 - z_2)
\;=\; \frac{dz_1 - dz_2}{z_1 - z_2}\,.
\end{equation}
This form has a simple pole along the diagonal $\Delta \subset
X \times X$ with residue~$1$.  The residue is intrinsic: it depends
on no coordinate, no metric, no level.  A double pole
$dz/(z-w)^2$ transforms under coordinate change and has no
coordinate-independent residue.  A regular $1$-form has residue
zero and extracts nothing from the collision.  The logarithmic
derivative is the unique $1$-form on $C_2(X)$ with a first-order
pole along~$\Delta$ and conformally invariant residue.

On $C_3(X)$, the three pullbacks
$\eta_{12}, \eta_{23}, \eta_{31}$ satisfy the \emph{Arnold relation}
\begin{equation}\label{eq:pf1-arnold}
\eta_{12} \wedge \eta_{23}
\;+\; \eta_{23} \wedge \eta_{31}
\;+\; \eta_{31} \wedge \eta_{12}
\;=\; 0\,.
\end{equation}
This restates, in differential forms, the partial-fractions identity
\[
\frac{1}{(z_1 - z_2)(z_2 - z_3)}
\;+\; \frac{1}{(z_2 - z_3)(z_3 - z_1)}
\;+\; \frac{1}{(z_3 - z_1)(z_1 - z_2)}
\;=\; 0\,,
\]
which holds because the left side, as a rational function of~$z_1$
with $z_2, z_3$ fixed, has simple poles at $z_2$ and $z_3$ whose
residues sum to zero.  Arnold proved that~\eqref{eq:pf1-arnold},
replicated across all triples
$\{i,j,k\} \subset \{1, \ldots, n\}$, generates the complete ideal
of relations in $H^*(\mathrm{Conf}_n(\mathbf{C}), \mathbf{C})$.

Under a coordinate change $z \mapsto w(z)$, the propagator
transforms as
$\eta_{ij} \mapsto \eta_{ij}
+ d\log\frac{w(z_i) - w(z_j)}{z_i - z_j}$.
Near the diagonal the correction becomes
$d\log w'(z_i)$: the standard cocycle for
$\mathrm{Diff}(X)$.
The propagator is a BRST ghost for diffeomorphisms: its failure to
be coordinate-invariant is precisely the failure that the Virasoro
algebra measures.


% ====================================================================
\subsection*{The bar complex}
% ====================================================================

A \emph{chiral algebra}~$\cA$ on~$X$ is a $\mathcal D_X$-module
(a sheaf with flat connection) equipped with a chiral bracket
\[
\mu^{\mathrm{ch}}
\;\colon\;
j_*j^*(\cA\boxtimes\cA)
\;\longrightarrow\;
\Delta_!(\cA)\,,
\]
where $j\colon X\times X\setminus\Delta
\hookrightarrow X\times X$ is the complement of the diagonal
and $\Delta_!\cA$ is the $!$-pushforward along the diagonal
embedding.  The source consists of sections of
$\cA \boxtimes \cA$ with poles of arbitrary order
along~$\Delta$; the chiral bracket extracts the singular part
and projects it onto the diagonal.  In local coordinates, the
chiral bracket encodes the operator product expansion:
\[
a(z)\,b(w)
\;\sim\;
\sum_{n\ge 0}\frac{a_{(n)}b(w)}{(z-w)^{n+1}}\,.
\]
The $n$-th products $a_{(n)}b\in\cA$ are the Fourier modes of
the bracket, extracted by
$a_{(n)}b = \operatorname{Res}_{z=w}
[a(z)\,b(w)\,(z{-}w)^n]$.
The zeroth product $a_{(0)}b$ is the Lie bracket
(first-order pole); the first product $a_{(1)}b$ carries the
symmetric/metric data (second-order pole).  Higher products
($n\ge 2$) are determined by the $\mathcal D_X$-module
structure.

Place sections $a_1,\dots,a_n\in\cA$ at
distinct points $z_1,\dots,z_n \in X$.  For a spanning tree~$T$
of the complete graph on $\{1,\dots,n\}$, form the decorated tensor
\[
\omega_T(a_1,\dots,a_n)
\;:=\;
a_1(z_1)\otimes\cdots\otimes a_n(z_n)
\;\otimes\;
\bigwedge_{(i,j)\in E(T)}\eta_{ij}\,.
\]
The form $\bigwedge_{(i,j)\in E(T)}\eta_{ij}$ is a section of
$\Omega^{n-1}_{\log}$ on the configuration space
$\mathrm{Conf}_n(X)$; its degree equals $|E(T)|=n-1$ since
$T$~is a tree.

The Fulton--MacPherson compactification
$\overline{\mathrm{FM}}_n(X)$ resolves $\mathrm{Conf}_n(X)$ by
iterated real oriented blowup along all diagonal strata.
Points of
$\overline{\mathrm{FM}}_n(X)\setminus\mathrm{Conf}_n(X)$
record the limiting direction and relative scale of approach as
points collide.  The boundary is a manifold with corners;
codimension-one faces $D_S$ are indexed by subsets
$S\subseteq\{1,\dots,n\}$ with $|S|\ge 2$, recording which
points have collided.  Each such face factors:
$D_S \simeq \overline{\mathrm{FM}}_{|S|}
\times \overline{\mathrm{FM}}_{n-|S|+1}$,
cluster times residual.

The \emph{ordered bar complex} is
\[
\barB^{\mathrm{ord}}_X(\cA)
\;=\;
\bigoplus_{n\ge 1}
\bigl(s^{-1}\bar\cA\bigr)^{\otimes n}
\;\otimes\;
\Gamma\bigl(\overline{\mathrm{FM}}_n(X),\;
\Omega^{n-1}_{\log}\bigr),
\]
with no $\Sigma_n$-quotient.  Here
$\bar\cA = \ker(\varepsilon\colon\cA\to\omega_X)$ is the
augmentation ideal, $s^{-1}$ denotes desuspension
($|s^{-1}a| = |a| - 1$), and the cohomological convention
$|d| = +1$ is used throughout.

The bar differential $d_{\barB}$ decomposes into three pieces:
\begin{equation}\label{eq:pf1-bar-diff}
d_{\barB}
\;=\;
d_{\mathrm{int}} + d_{\mathrm{res}} + d_{\mathrm{dR}}\,.
\end{equation}
The three components: $d_{\mathrm{int}}$ is the internal
differential of~$\cA$ (acting on each tensor factor);
$d_{\mathrm{dR}}$ is the de~Rham differential on logarithmic
forms on~$\overline{\mathrm{FM}}_n(X)$; and $d_{\mathrm{res}}$
is the residue differential, which extracts OPE data at collision
divisors.  For a codimension-one face $D_{\{i,j\}}$
corresponding to the collision $z_i \to z_j$:
\begin{equation}\label{eq:pf1-res-diff}
d_{\mathrm{res}}^{(ij)}
\bigl(
a_1\otimes\cdots\otimes a_n\otimes\eta_{ij}\wedge\omega'
\bigr)
\;=\;
\pm\,
a_1\otimes\cdots\otimes
\widehat{a_i}\otimes\cdots\otimes
\widehat{a_j}\otimes\cdots\otimes
a_{i(p)}a_j
\;\otimes\;
\omega',
\end{equation}
where $p$ is the pole order read from the form-theoretic coefficient,
hats denote omission, and $a_{(p)}b$ is the $p$-th Fourier mode of
the OPE\@.  The bar differential does not merely detect singularities;
it extracts them, mode by mode, through the logarithmic kernel.

\subsection*{Why $d^2 = 0$}
% ====================================================================

Expand $(d_{\mathrm{int}}+d_{\mathrm{res}}+d_{\mathrm{dR}})^2$
into nine terms.  Three vanish on the diagonal:
$d_{\mathrm{int}}^2 = 0$ ($\cA$ is a cochain complex),
$d_{\mathrm{dR}}^2 = 0$ (de~Rham), and
$d_{\mathrm{int}}\,d_{\mathrm{dR}}
+ d_{\mathrm{dR}}\,d_{\mathrm{int}} = 0$
(Leibniz rule).
The remaining terms pair into three cancellation mechanisms,
governed by index overlap:

\begin{enumerate}[label=(\alph*)]
\item \emph{Disjoint supports.}\;
When $\{i,j\}\cap\{k,l\}=\emptyset$, the faces
$D_{\{i,j\}}$ and $D_{\{k,l\}}$ are transverse in
$\overline{\mathrm{FM}}_n(X)$.  The residue operators commute:
$d_{\mathrm{res}}^{(ij)}\,d_{\mathrm{res}}^{(kl)}
+ d_{\mathrm{res}}^{(kl)}\,d_{\mathrm{res}}^{(ij)} = 0$.

\item \emph{Shared index.}\;
When $\{i,j\}\cap\{k,l\} = \{j=k\}$, the triple collision
$z_i = z_j = z_l$ sits in a codimension-two stratum.  The
form-theoretic contribution involves
$\eta_{ij}\wedge\eta_{jl}$.  The Arnold
relation~\eqref{eq:pf1-arnold} identifies this with a cyclic
sum over the three ordered approaches to the triple-collision
point.  On the algebraic side, the three terms produce
$(a_{i(m)}a_j)_{(n)}a_l$,\; $(a_{j(m)}a_l)_{(n)}a_i$,\;
$(a_{l(m)}a_i)_{(n)}a_j$ (with appropriate pole orders),
whose sum vanishes by the \emph{Borcherds identity}:
\begin{equation}\label{eq:pf1-borcherds}
\boxed{
\sum_{j\ge 0}\binom{m}{j}(a_{(n+j)}b)_{(m+k-j)}c
\;=\;
\sum_{j\ge 0}(-1)^j\binom{n}{j}
\Bigl(
a_{(m+n-j)}(b_{(k+j)}c)
-(-1)^n b_{(n+k-j)}(a_{(m+j)}c)
\Bigr).
}
\end{equation}
The Arnold relation kills the form; the Borcherds identity
kills the algebra.  Both are needed, and both are consumed.

\item \emph{Same pair.}\;
$d_{\mathrm{res}}^{(ij)}\,d_{\mathrm{res}}^{(ij)} = 0$ by the
sign rule on desuspensions: extracting a residue twice at the
same diagonal vanishes for degree reasons.
\end{enumerate}

The result: $d_{\barB}^2 = 0$ if and only if the $n$-th products
satisfy the Borcherds identity~\eqref{eq:pf1-borcherds}.  At $n=0$
the Borcherds identity is the Jacobi identity for the Lie
bracket~$a_{(0)}$; at $n=1$, $m=k=0$, the derivation property
of~$a_{(0)}$ with respect to~$a_{(1)}b$.  The Borcherds identity is
the complete axiom system for a chiral algebra; the bar construction
accepts exactly chiral algebras.


% ====================================================================
\subsection*{Two coproducts}
% ====================================================================

The ordered bar $\barB^{\mathrm{ord}}(\cA)$ carries the
\emph{deconcatenation coproduct}:
\begin{equation}\label{eq:pf1-deconc}
\Delta[s^{-1}a_1|\cdots|s^{-1}a_n]
\;=\;
\sum_{i=0}^{n}
[s^{-1}a_1|\cdots|s^{-1}a_i]
\;\otimes\;
[s^{-1}a_{i+1}|\cdots|s^{-1}a_n]\,,
\end{equation}
splitting an ordered sequence at every position: $n+1$ terms.
This is coassociative but not cocommutative, an
$E_1$-coalgebra reflecting the linear ordering.

The $\Sigma_n$-coinvariant quotient
\[
\barB^{\Sigma}_X(\cA)
\;=\;
\bigoplus_{n\ge 1}
\bigl(s^{-1}\bar\cA\bigr)^{\otimes n}_{\Sigma_n}
\;\otimes\;
\Gamma\bigl(\overline{\mathrm{FM}}_n(X),\;
\Omega^{n-1}_{\log}\bigr)
\]
carries the \emph{coshuffle coproduct}, summing over all $2^n$
bipartitions of an unordered set: coassociative and cocommutative,
an $E_\infty$-coalgebra, the factorization coalgebra of Beilinson
and Drinfeld on~$\operatorname{Ran}(X)$.

The ordered bar is the natural primitive.  The symmetric bar is a
quotient that loses data: the non-cocommutative structure of the
ordered bar carries the spectral $R$-matrix, the KZ associator, and
the full quantum group deformation, none of which survive the
passage to $\Sigma_n$-coinvariants.  How much is lost becomes
visible in the first examples.


% ====================================================================
\subsection*{Logarithm, exponential, dual}
% ====================================================================

The bar construction is a categorical logarithm:
$B(\cA\otimes\cA') \simeq B(\cA)\oplus B(\cA')$.
The logarithmic form $\eta_{ij}=d\log(z_i-z_j)$ is the integral
kernel; the Arnold relation is the cocycle condition for~$\log$.

The cobar functor $\Omega$ inverts the bar on the Koszul locus:
$\Omega(\barB(\cA))\xrightarrow{\sim}\cA$ (Theorem~B).
This is bar-cobar \emph{inversion}: the cobar recovers the
\emph{original} algebra, not a dual.

A separate operation produces the dual.  Take bar cohomology
$\cA^{\mathrm i} := H^*(B(\cA))$, a graded coalgebra, and apply
graded linear duality: $\cA^! := (\cA^{\mathrm i})^\vee$.
This is the \emph{Koszul dual algebra}.  Verdier duality
on $\mathrm{Ran}(X)$ produces the richer chain-level
object~$\cA^!_\infty$:
\begin{equation}\label{eq:pf1-verdier}
\mathbb D_{\mathrm{Ran}}\,\barB_X(\cA)
\;\simeq\;
\cA^!_\infty.
\end{equation}
Verdier duality converts coalgebra to algebra, so $\cA^!_\infty$
is a factorization \emph{algebra}, retaining full chain-level data.
That its cohomology recovers~$\cA^!$ on the Koszul locus is the
content of Theorem~A, not a tautology.

Three functors act on $B(\cA)$ and produce three distinct outputs:
\[
\begin{array}{lcll}
\Omega(B(\cA)) &\simeq& \cA
  & \text{(inversion: recovers the original)},\\[3pt]
\mathbb D_{\mathrm{Ran}}(B(\cA)) &\simeq& \cA^!_\infty
  & \text{(Verdier: the homotopy dual)},\\[3pt]
(H^*(B(\cA)))^\vee &=& \cA^!
  & \text{(linear duality: the strict dual)}.
\end{array}
\]
Confusing any two is a category error.

When~$\cA$ is chirally Koszul, $H^*(B(\cA))$ concentrates in bar
degree~$1$, the bar-to-cobar spectral sequence collapses at~$E_2$,
and $\cA^!$ governs a linear resolution of~$\cA$.  Classical Koszul
duality (Priddy, Beilinson--Ginzburg--Soergel, Loday--Vallette) is
recovered on $X=\mathbb A^1$: the affine line is contractible, the
FM cooperad becomes quasi-isomorphic to the associative cooperad,
and the bar complex reduces to the classical Chevalley--Eilenberg
complex.


% ====================================================================
\subsection*{The Heisenberg algebra}
% ====================================================================

The Heisenberg chiral algebra $\mathcal H_k$ is generated by a
single field~$\alpha$ of conformal weight~$1$ with OPE
\[
\alpha(z)\,\alpha(w)
\;\sim\;
\frac{k}{(z-w)^2}\,.
\]
Only the second-order pole is present: $\alpha_{(1)}\alpha
= k\cdot\mathbf 1$ and $\alpha_{(n)}\alpha = 0$ for $n\neq 1$.
No first-order pole: $\alpha_{(0)}\alpha = 0$, no Lie bracket.
All triple products vanish since
$\mathbf 1_{(n)} = 0$ for $n\ge 0$.

The bar complex at tensor degree~$n$ is
\[
\barB^n_X(\mathcal H_k)
\;=\;
(s^{-1}\alpha)^{\otimes n}
\;\otimes\;
\Gamma\bigl(\overline{\mathrm{FM}}_n(X),\;
\Omega^{n-1}_{\log}\bigr).
\]
At tensor degree~$1$: a one-dimensional space, $d_{\barB} = 0$.
At tensor degree~$2$: the bar differential extracts the full
chiral product:
$d_{\mathrm{res}}(s^{-1}\alpha\otimes s^{-1}\alpha
\otimes\eta_{12}) = \alpha_{(1)}\alpha = k\cdot\mathbf{1}$.
The level~$k$ is visible already at genus~$0$.  At tensor
degree~$n \ge 3$: every residue at a collision stratum
$D_{\{i,j\}}$ produces $\alpha_{(1)}\alpha
= k\cdot\mathbf 1$ on the pair~$(i,j)$; since
$\mathbf 1_{(p)}\alpha = 0$ for $p\ge 0$, no further
contractions occur.  All structure is captured at tensor
degree~$2$: the bar complex is quadratic.
The Heisenberg algebra is chirally Koszul.

The modular characteristic:
$\kappa(\mathcal H_k) = k$.
The Koszul dual:
\[
\mathcal H_k^!
\;=\;
\bigl(\mathrm{Sym}^{\mathrm{ch}}(V^*),\;
m_0 = -k\omega\bigr),
\]
where $V^*$ is the linear dual of the generating space and
$m_0 = -k\omega$ is the curvature, satisfying the curved
$A_\infty$ relation $m_1^{\,2}(a) = [m_0, a]$ (commutator,
with the minus sign).  The differential $m_1$ does not square to
zero; its failure is controlled by the curvature.

\medskip
\noindent\textbf{The ordered bar and its $R$-matrix.}\enspace
The ordered bar complex $\barB^{\mathrm{ord}}(\mathcal H_k)$
carries the collision residue
\begin{equation}\label{eq:pf1-heis-rmatrix}
r(z) \;=\; \frac{k}{z}\,,
\end{equation}
a scalar rational function with a single simple pole.  The
second-order OPE pole $k/(z-w)^2$ drops by one order under the
logarithmic kernel $d\log(z-w)$: the bar differential absorbs one
power of $(z-w)$, producing the first-order $r$-matrix pole~$k/z$.
Applying the averaging map:
\[
\mathrm{av}\bigl(r(z)\bigr) \;=\; k
\;=\; \kappa(\mathcal H_k).
\]
For the Heisenberg algebra, the $R$-matrix \emph{is} the modular
characteristic: $r(z) = k/z$ is already $\Sigma_2$-invariant
(scalar, not matrix), so averaging loses nothing.  The two Koszul
dualities coincide, the ordered bar carries no more information than
the symmetric bar, and the entire structure collapses to a
single number~$k$.

The deficiency: no first-order pole, no Lie bracket, no matrix
structure in the collision residue, no quantum group.  Everything
that makes representation theory nontrivial is absent.  What
Chriss--Ginzburg demands is the simplest example where the $R$-matrix
is genuinely matrix-valued: the first algebra with a first-order
pole.


% ====================================================================
\subsection*{The Kac--Moody OPE}
% ====================================================================

The affine Kac--Moody algebra $\widehat{\mathfrak g}_k$ at
level~$k$ is generated by currents
$J^a$, $a = 1,\dots,\dim\mathfrak g$, with OPE
\begin{equation}\label{eq:pf1-km-ope}
J^a(z)\,J^b(w)
\;\sim\;
\frac{f^{ab}{}_{c}\,J^c(w)}{z-w}
\;+\;
\frac{k\,\kappa^{ab}}{(z-w)^2}\,,
\end{equation}
where $f^{ab}{}_{c}$ are the structure constants of the simple Lie
algebra~$\mathfrak g$ and $\kappa^{ab}$ is the Killing form
normalised so that long roots have squared length~$2$.  Both poles
are present.  The first-order pole carries the Lie bracket
$J^a_{(0)}J^b = f^{ab}{}_{c}J^c$; the second-order pole carries
the invariant inner product
$J^a_{(1)}J^b = k\,\kappa^{ab}\cdot\mathbf 1$.

The bar complex of $\widehat{\mathfrak g}_k$ has nontrivial
contributions at all tensor degrees: the first-order pole generates
structure constants from tree-level OPE contractions, and the
second-order pole contributes the metric data.  The chiral bracket
decomposes:
\begin{equation}\label{eq:pf1-km-bracket}
\mu^{\mathrm{ch}}
\;=\;
\underbrace{f^{ab}{}_{c}\,J^c\cdot\eta_{12}}_{\text{Lie bracket
(first-order pole)}}
\;+\;
\underbrace{k\,\kappa^{ab}\cdot\mathbf 1\cdot
d\eta_{12}}_{\text{inner product (second-order pole)}}.
\end{equation}
The Lie bracket drives the tree-level bar differential
(the combinatorics of the Chevalley--Eilenberg complex); the inner
product supplies the curvature~$\kappa$.  The coexistence of both
poles, absent in the Heisenberg case, makes
$\widehat{\mathfrak g}_k$ the first nontrivial example.

\medskip
\noindent\textbf{The Casimir $R$-matrix.}\enspace
The ordered bar complex of $\widehat{\mathfrak g}_k$ carries the
collision residue
\begin{equation}\label{eq:pf1-km-rmatrix}
r(z) \;=\; \frac{\Omega}{z}\,,
\qquad
\Omega \;=\; \sum_a J^a \otimes J_a\,,
\end{equation}
where $\Omega \in \mathfrak g \otimes \mathfrak g$ is the Casimir
tensor.  This is matrix-valued: $\Omega$ carries the full
Lie-algebraic structure that was invisible in the Heisenberg
scalar~$k/z$.  The averaging map collapses the Casimir to a
scalar:
\begin{equation}\label{eq:pf1-km-kappa}
\mathrm{av}\bigl(k\Omega/z\bigr)
\;=\;
\frac{\dim(\mathfrak g)\cdot(k+h^\vee)}{2h^\vee}
\;=\;
\kappa(\widehat{\mathfrak g}_k).
\end{equation}
The matrix structure of~$\Omega$, erased by
$\Sigma_2$-coinvariance, encodes the representation theory
of~$\mathfrak g$.  At this point the ordered bar has strictly more
information than the symmetric bar: the passage from
$r(z) = k\Omega/z$ to $\kappa(\widehat{\mathfrak g}_k)$ discards
the Casimir tensor, which is precisely the datum that builds the
Yangian~$Y(\mathfrak g)$.

The Yangian is the Koszul dual of the ordered bar: the quantum group
that deforms the enveloping algebra $U(\mathfrak g[z])$ of
polynomial currents.  The symmetric bar produces the chiral Koszul
dual $\widehat{\mathfrak g}_k^!$, the Verdier-dual factorization
algebra on~$\operatorname{Ran}(X)$.  The two dualities coexist
because the bar complex carries two operadic colours.

The five main theorems of this monograph are the invariants of
the ordered Maurer--Cartan element~$\Theta^{E_1}_\cA$ that survive
the averaging map to~$\Sigma_n$-coinvariants.
The adjunction and Verdier intertwining (Theorem~A), the bar-cobar
inversion (Theorem~B), the complementarity decomposition
(Theorem~C), the modular characteristic~$\kappa$ (Theorem~D), and
the chiral Hochschild polynomial growth (Theorem~H) are all
coinvariant projections of the richer $E_1$-data.
The ordered story is the primitive; the modular story is its shadow.


% ====================================================================
\subsection*{Genus $0$ and genus $\ge 1$}
% ====================================================================

Everything above takes place at genus~$0$: the propagator
$\eta_{12} = d\log(z_1-z_2)$ is globally defined
on~$\mathbb P^1$, and the Arnold relation holds without correction.
On a curve of genus $g \geq 1$, the function $z_1 - z_2$ has no
global meaning: the curve is compact, and any meromorphic function
with a single simple zero must also have a simple pole.

The replacement is the \emph{prime form} $E(z_1,z_2)$, a section of
$K^{-1/2}\boxtimes K^{-1/2}$ (not $K^{+1/2}\boxtimes K^{+1/2}$),
vanishing to first order along the diagonal with no other zeros.
The genus-$g$ propagator
$\eta^{(g)}_{12} = d\log E(z_1,z_2) + \text{(period correction)}$
is single-valued but no longer holomorphic; the period correction
couples to the Hodge bundle
$\mathbb E \to \overline{\mathcal M}_g$ through the period
matrix~$\Omega$.  The propagator $d\log E(z_1,z_2)$ has conformal
weight~$1$ in both variables, regardless of the conformal weight of
the field being sewed: all edges of the genus-$g$ graph sum use the
standard Hodge bundle $\mathbb E_1 = R^0\pi_*\omega_\pi$.

The Arnold relation acquires a defect, and the fibrewise bar
differential satisfies
\begin{equation}\label{eq:pf1-curvature}
d_{\mathrm{fib}}^{\,2}
\;=\;
\kappa(\cA)\cdot\omega_g\,,
\end{equation}
where $\omega_g$ is the Arakelov form on the fibre of
$\pi\colon\mathcal C_g \to \overline{\mathcal M}_g$ and
$\kappa(\cA)\in\Bbbk$ is the modular characteristic,
determined entirely by the leading OPE singularity.  Explicit
values: $\kappa(\mathcal H_k) = k$,\;
$\kappa(\widehat{\mathfrak g}_k)
= \dim(\mathfrak g)\cdot(k+h^\vee)/(2h^\vee)$,\;
$\kappa(\mathrm{Vir}_c) = c/2$.
Period integrals restore nilpotence: the Gauss--Manin connection
absorbs the fibrewise curvature, and the total differential
satisfies $D_g^2 = 0$.

Genus~$0$ is uncurved algebra.  Genus $\geq 1$ is curved
geometry.  The entire genus tower is controlled by a single
scalar: the curvature of the bar complex on a family of curves.

\begin{remark}[The zero-dimensional shadow]\label{rem:pf1-zero-dim}
When $X=\operatorname{Spec}\Bbbk$ is a point, configuration spaces
collapse, logarithmic forms vanish, and the chiral bar complex
recovers the classical bar construction of an associative
algebra~$A$.  A quadratic algebra $A = T(V)/(R)$ has quadratic dual
$A^! = T(V^*)/(R^\perp)$; when~$A$ is Koszul, the bar cohomology
concentrates in diagonal degrees and
$H^*(B(A)) \cong (A^!)^\vee$.  The operadic instances
$\mathrm{Com}^! = \mathrm{Lie}$ and
$\mathrm{Sym}^! = \Lambda$ recover the Chevalley--Eilenberg
complex as a restricted bar construction.
\end{remark}
